\documentclass[11pt]{article}
\usepackage[T1]{fontenc}
\usepackage{lmodern,amsmath,amssymb,amsthm,microtype}
\usepackage[margin=1in]{geometry}
\usepackage[hidelinks]{hyperref}
\hypersetup{pdftitle={Frobenius splitting and singular loci of prime-size permanents},pdfauthor={Ying Xie},pdfcreator={},pdfproducer={}}
\ifdefined\pdfinfoomitdate\pdfinfoomitdate=1\fi
\ifdefined\pdfsuppressptexinfo\pdfsuppressptexinfo=15\fi
\ifdefined\pdftrailerid\pdftrailerid{}\fi
\newtheorem{theorem}{Theorem}[section]
\newtheorem{proposition}[theorem]{Proposition}
\newtheorem{lemma}[theorem]{Lemma}
\newtheorem{corollary}[theorem]{Corollary}
\theoremstyle{definition}\newtheorem{definition}[theorem]{Definition}
\theoremstyle{remark}\newtheorem{remark}[theorem]{Remark}
\DeclareMathOperator{\per}{per}
\DeclareMathOperator{\Crit}{Crit}
\DeclareMathOperator{\Sing}{Sing}
\DeclareMathOperator{\codim}{codim}
\DeclareMathOperator{\rank}{rank}
\DeclareMathOperator{\htop}{ht}
\DeclareMathOperator{\str}{str}
\newcommand{\F}{\mathbb F}
\newcommand{\A}{\mathbb A}
\newcommand{\Q}{\mathbb Q}
\newcommand{\Z}{\mathbb Z}
\newcommand{\m}{\mathfrak m}
\newcommand{\pp}{\mathfrak p}
\numberwithin{equation}{section}
\title{Frobenius splitting and singular loci\\of prime-size permanents}
\author{Ying Xie\\Kennesaw State University\\\texttt{yxie2@kennesaw.edu}}
\date{}
\begin{document}
\maketitle
\begin{abstract}
For each prime $p$, we compute a Frobenius trace coefficient associated
with the maximal permanents of a generic $(p-1)\times p$ matrix.
A cyclic permutation of the columns makes this coefficient equal to
one in characteristic $p$. We prove that the maximal permanents form
a radical complete intersection of height $p$ in characteristic $p$
and in characteristic zero. A local differential argument bounds every
Jacobian-rank stratum and gives ambient codimension exactly $2p$ for
the singular locus of the $p\times p$ permanent in both
characteristics. A dimension estimate obtained by fixing a nonzero
border entry propagates the bound to arbitrary matrix sizes. In
characteristic zero, if $q$ is the largest prime at most $n$, the
singular codimension lies between $2q$ and $2n$, and hence is
asymptotic to $2n$. We also obtain exact critical codimensions for
partial row expansions and exact strength at prime sizes.
\end{abstract}

\section{Introduction and main results}
Let $X=(x_{ij})$ be a generic $n\times n$ matrix over a field $k$.
Its permanent is
\[
 \per_n(X)=\sum_{\sigma\in S_n}\prod_{i=1}^n x_{i,\sigma(i)}.
\]
The first partial derivatives of $\per_n$ are its permanental
cofactors. Thus the geometry of the singular locus of the permanent
is closely related to the codimension of rectangular permanental
ideals.

For an $r\times n$ generic matrix $A$, with $r\le n$, let $P_r(A)$
denote the ideal of its maximal permanents. Boralevi, Carlini,
Micha\l{}ek, and Ventura conjecture that the maximal permanents of
a generic $r\times(r+1)$ matrix form a complete intersection
\cite[Conjecture~3.4]{BCMV}. Their work establishes several small
cases and gives a uniform lower bound of six for the ambient singular
codimension of $\per_n$ when $n\ge6$
\cite[Theorem~4.23]{BCMV}. We treat all row sizes one less than a
prime, and obtain asymptotically sharp bounds for square permanents
of arbitrary size.

Throughout, dimensions are geometric dimensions: they may be
computed after extension to an algebraic closure. The critical locus
of a polynomial is the common zero set of all its first derivatives.
For the permanent it lies on the hypersurface, in every
characteristic, by expansion along a row. We write
\[
 s_k(n)=\codim_{\A^{n^2}_k}\Crit(\per_n),\qquad
 c_k(r,n)=\htop P_r(A).
\]
All codimensions in this paper are measured in matrix space, rather
than inside the permanental hypersurface.

\begin{theorem}\label{thm:prime}
Let $p$ be prime, and let $A$ be a generic $(p-1)\times p$ matrix.
Over every field of characteristic $p$ or zero, the ideal $P_{p-1}(A)$
is a geometrically reduced complete intersection of height $p$.
Its Hilbert series is
\begin{equation}\label{eq:hilbert}
 \frac{(1-z^{p-1})^p}{(1-z)^{p(p-1)}}.
\end{equation}
In characteristic $p$, its quotient is Frobenius split. In these
same two characteristics,
\begin{equation}\label{eq:primecodim}
 s_k(p)=2p.
\end{equation}
\end{theorem}

The splitting in this theorem uses the prime equal to the number of
columns. The assertion is not a statement of splitting in every
positive characteristic for a fixed matrix size. Characteristic
dependence of permanental $F$-singularities is studied in \cite{Chau}.

\begin{theorem}\label{thm:alln}
Suppose $k$ has characteristic zero and $n\ge2$. If $q$ is the
largest prime at most $n$, then
\begin{equation}\label{eq:alln}
 2q\le s_k(n)\le2n,\qquad q\le c_k(n-1,n)\le n.
\end{equation}
Consequently,
\[
 \lim_{n\to\infty}\frac{s_k(n)}{2n}=1,
 \qquad
 \lim_{n\to\infty}\frac{c_k(n-1,n)}{n}=1.
\]
\end{theorem}

The proof begins with a coefficient identity. If $f_j$ is the
maximal permanent of $A$ omitting column $j$ and $H=\prod_jf_j$,
then
\begin{equation}\label{eq:introtrace}
 \left[\prod_{i=1}^{p-1}\prod_{j=1}^{p}a_{ij}^{p-1}\right]
 H^{p-1}=1\quad\text{in }\F_p.
\end{equation}
This is a specific instance of the relation between finite-field
point counting and Frobenius splitting developed by Knutson
\cite{Knutson}. We give two direct cyclic proofs. We then derive a
local rank estimate from the nonzero coefficient, without assuming a
complete intersection in advance. This separates the dimension
argument from the later splitting and reducedness assertions.

\begin{theorem}\label{thm:partial}
Use the maximal permanents $f_1,\ldots,f_p$ above, and let
$S\subseteq\{1,\ldots,p\}$ have cardinality $s\ge1$. Over a field
of characteristic $p$ or zero, put
\[
 g_S(A,u)=\sum_{j\in S}u_jf_j(A).
\]
In the affine space with coordinates $A$ and $(u_j)_{j\in S}$,
\[
 \codim\Crit(g_S)=2s.
\]
Its strength, counted as the minimum number of products of
positive-degree homogeneous polynomials, is exactly $s$.
\end{theorem}

The strength consequence for arbitrary square size is
\begin{equation}\label{eq:introstrength}
 q\le\str(\per_n)\le n\qquad(\operatorname{char}k=0).
\end{equation}
This convention for strength counts products, with no subtraction of
one. Neither this invariant nor the singular codimension alone
establishes an unrestricted arithmetic-circuit lower bound.

\section{The cyclic trace coefficient}
Let $N=p(p-1)$, let $A=(a_{ij})$ be a generic $(p-1)\times p$
matrix, and set
\[
 f_j=\per(A_{\widehat j}),\qquad
 H=\prod_{j=1}^{p}f_j,\qquad M=\prod_{i,j}a_{ij}.
\]
Both $H$ and $M$ have degree $N$.

\begin{proposition}\label{prop:cyclic}
For every prime $p$, the integer coefficient $[M^{p-1}]H^{p-1}$
is congruent to one modulo $p$.
\end{proposition}
\begin{proof}
For a nonnegative integer $e$,
\[
 \sum_{a\in\F_p}a^e=
 \begin{cases}-1,&e>0\text{ and }p-1\mid e,\\
 0,&\text{otherwise}.
 \end{cases}
\]
In a monomial of total degree $N(p-1)$, a nonzero sum over
$\F_p^N$ requires every exponent to be a positive multiple of
$p-1$. Thus each exponent must be exactly $p-1$. Since $N$ is
even, we obtain
\begin{equation}\label{eq:pointsum}
 [M^{p-1}]H^{p-1}=\sum_{A\in\operatorname{Mat}_{p-1,p}(\F_p)}H(A)^{p-1}
 \quad\text{in }\F_p.
\end{equation}
Cyclic permutation of the $p$ columns preserves the summand.
Every orbit has size one or $p$, so only fixed matrices contribute
modulo $p$. A fixed matrix has all columns equal to a vector
$a=(a_1,\ldots,a_{p-1})^{\mathsf t}$. Its maximal permanents all
equal $(p-1)!\prod_i a_i$. Hence its contribution is one exactly
when every $a_i$ is nonzero, and is zero otherwise. The sum is
therefore $(p-1)^{p-1}=1$ in $\F_p$. The argument also applies when
$p=2$.
\end{proof}

There is also a coefficient-level proof, which makes the expansion
multiplicities explicit. It will be useful as an independent check
on the point-counting argument.

\begin{proof}[A second proof of Proposition~\ref{prop:cyclic}]
Index columns by $\F_p$ and label the $p-1$ occurrences of each
$f_j$ by $h\in\{1,\ldots,p-1\}$. A term in the product is a
family of bijections
\[
 \sigma_{j,h}:\{1,\ldots,p-1\}\longrightarrow\F_p\setminus\{j\}.
\]
It contributes to the required coefficient if, for each row $i$ and
column $a$, exactly $p-1$ pairs $(j,h)$ satisfy
$\sigma_{j,h}(i)=a$. Call this the balance condition. The transformation
\[
 (T\sigma)_{j,h}(i)=\sigma_{j-1,h}(i)+1
\]
preserves balance and has order $p$. A fixed family is precisely
one of the form $\sigma_{j,h}(i)=\tau_h(i)+j$, where each $\tau_h$
is a bijection from $\{1,\ldots,p-1\}$ to $\F_p^*$.
Every such family is balanced, since varying $j$ visits each column
once for fixed $i,h$. There are $((p-1)!)^{p-1}$ fixed families.
Wilson's congruence shows that this count is one modulo $p$,
including the immediate case $p=2$. All other orbits have size $p$.
\end{proof}

\section{A local consequence of a nonzero coefficient}
The next statements apply to a general product of homogeneous
polynomials whose degrees sum to the number of variables.

\subsection{Differential nonmembership}
Let $R=k[x_1,\ldots,x_N]$ in characteristic $p$, and let
$f_1,\ldots,f_t$ be positive-degree homogeneous polynomials with
$\sum_i\deg f_i=N$. Put $H=\prod_i f_i$ and assume
\begin{equation}\label{eq:nonzero}
 c=[x_1^{p-1}\cdots x_N^{p-1}]H^{p-1}\ne0.
\end{equation}
For an ideal $J$, write $J^{[p]}=(a^p:a\in J)$.

\begin{lemma}\label{lem:diff}
For every prime $\pp$ of $R$, writing
$\m=\pp R_{\pp}$, one has $H^{p-1}\notin\m^{[p]}$.
The same statement holds for the product of any subset of the $f_i$,
raised to the power $p-1$.
\end{lemma}
\begin{proof}
Consider the composition of ordinary differential operators
\[
 D=\prod_{i=1}^{N}\left(\frac{\partial}{\partial x_i}\right)^{p-1}.
\]
Any monomial with an exponent below $p-1$ is annihilated, and the
total degree $N(p-1)$ leaves only the indicated monomial. Therefore
\begin{equation}\label{eq:D}
 D(H^{p-1})=((p-1)!)^N c\in k^*.
\end{equation}
Each derivation extends to $R_{\pp}$ and preserves $\m^{[p]}$:
for $a\in\m$ and $b\in R_{\pp}$,
$\partial(a^pb)=a^p\partial b$ in characteristic $p$.
This also holds for fractions. Thus their composition preserves
$\m^{[p]}$, and \eqref{eq:D} rules out membership of $H^{p-1}$.
If a subproduct belonged to $\m^{[p]}$, multiplying by the remaining
factors would give the forbidden membership for $H^{p-1}$.
\end{proof}

\subsection{Two elementary local identities}
\begin{lemma}\label{lem:pigeonhole}
If $(S,\m)$ is a regular local ring of dimension $h$ and
characteristic $p$, then
\begin{equation}\label{eq:pigeonhole}
 \m^{h(p-1)+1}\subseteq\m^{[p]}.
\end{equation}
\end{lemma}
\begin{proof}
The maximal ideal has $h$ generators. In every monomial of degree
$h(p-1)+1$ in those generators, one exponent is at least $p$.
\end{proof}

\begin{lemma}\label{lem:change}
Let $S$ be a local ring of characteristic $p$, let
$I=(h_1,\ldots,h_s)$, and let $g=Uh$ for $U\in\operatorname{GL}_s(S)$.
Then
\begin{equation}\label{eq:change}
 \prod_{i=1}^{s}g_i^{p-1}
 \equiv\det(U)^{p-1}\prod_{i=1}^{s}h_i^{p-1}\pmod{I^{[p]}}.
\end{equation}
\end{lemma}
\begin{proof}
Work first in $S[y_1,\ldots,y_s]/(y_1^p,\ldots,y_s^p)$.
Its top homogeneous component is free of rank one, generated by
$\prod_i y_i^{p-1}$. A scaling $y_i\mapsto ay_i$ by a unit acts
on this component by $a^{p-1}$. A permutation acts trivially, in
agreement with its determinant to the power $p-1$, also when $p=2$.
An elementary shear $y_i\mapsto y_i+ay_j$ fixes the component:
every term other than $y_i^{p-1}$ in the expansion has exponent at
least $p$ on $y_j$ after multiplication by $y_j^{p-1}$.
These operations generate the invertible matrices over a local
ring, by elimination using a unit entry. They therefore prove the
determinant character asserted in \eqref{eq:change}. Substitution
$y_i=h_i$ gives the result. Entries of $U$ need not be constants.
\end{proof}

\subsection{Height and Jacobian-rank strata}
Assume now that $k$ is perfect. Choose a subset of $s$ of the
polynomials in \eqref{eq:nonzero}, relabel it $h_1,\ldots,h_s$, and
put $I_S=(h_1,\ldots,h_s)$. Let $J_h$ be their $s\times N$
Jacobian matrix. For $0\le r\le s$, define the locally closed set
\[
 Z_r=\{x\in V(I_S):\rank J_h(x)=r\}.
\]

\begin{proposition}\label{prop:rank}
Under \eqref{eq:nonzero}, one has
\begin{equation}\label{eq:rankstrata}
 \htop I_S=s,\qquad \codim_{\A^N} Z_r\ge2s-r
\end{equation}
for every nonempty rank stratum. The generators of $I_S$ form a
regular sequence.
\end{proposition}
\begin{proof}
At any prime $\pp\supseteq I_S$, set $S_0=R_{\pp}$,
$\m=\pp S_0$, and $h=\htop\pp$. If $h<s$, then
\[
 \prod_{i=1}^{s}h_i^{p-1}\in\m^{s(p-1)}
 \subseteq\m^{h(p-1)+1}\subseteq\m^{[p]},
\]
contradicting Lemma~\ref{lem:diff}. The height theorem gives the
opposite bound, so $\htop I_S=s$. Cohen--Macaulayness of the
polynomial ring gives the regular-sequence assertion.

For the refined estimate, take the generic point $\pp$ of the
closure of an irreducible component of $Z_r$. Let
$K=\operatorname{Frac}(R/\pp)$ and $h=\htop\pp$.
The span of the classes of the $h_i$ in $\m/\m^2$ has dimension
$r$. To justify this point, use the conormal sequence
\begin{equation}\label{eq:conormal}
 \m/\m^2\longrightarrow\Omega_{R/k}\otimes_R K
       \longrightarrow\Omega_{K/k}\longrightarrow0.
\end{equation}
The first space has dimension $h$ because $R_{\pp}$ is regular;
the middle one has dimension $N$. Since $k$ is perfect and $K/k$
is a finitely generated field extension,
$\dim_K\Omega_{K/k}=\operatorname{trdeg}_kK=N-h$.
Exactness of \eqref{eq:conormal} therefore makes its first map
injective. Its values on the generator classes are exactly the
Jacobian rows evaluated over $K$.

Row reduction over $K$, lifted to an invertible matrix over
$R_{\pp}$, replaces the $h_i$ by $g_i$ with
\[
 g_1,\ldots,g_r\in\m,\qquad
 g_{r+1},\ldots,g_s\in\m^2.
\]
Lemma~\ref{lem:change}, together with $I_S^{[p]}\subseteq\m^{[p]}$,
shows that $\prod_i g_i^{p-1}\notin\m^{[p]}$.
But this product belongs to $\m^{(2s-r)(p-1)}$.
If $h<2s-r$, Lemma~\ref{lem:pigeonhole} would force membership
in $\m^{[p]}$, a contradiction. Thus $h\ge2s-r$, as required.
\end{proof}

The rank estimate did not use reducedness or a splitting on the
quotient. Its only coefficient input was \eqref{eq:nonzero}. We
will use this separation when applying it to critical loci.

\section{Compatible splitting and characteristic-zero transfer}
\subsection{The prime characteristic}
Apply Proposition~\ref{prop:cyclic} to the maximal permanents
$f_1,\ldots,f_p$. Proposition~\ref{prop:rank} proves the height
and regular-sequence statements over an algebraic closure of $\F_p$,
and hence over every field of characteristic $p$ by faithfully flat
descent. We next establish reducedness and the claimed splitting.

Choose an additive coefficient map $\lambda:k\to k$ with
$\lambda(a^pb)=a\lambda(b)$ and $\lambda(1)=1$. It exists over
any field: project $k$ onto its $1$-coordinate in a $k^p$-basis
containing $1$, and then take the $p$th root. For the $N$ entries
of $A$, define
\[
 T(bx^{\alpha})=
 \begin{cases}
 \lambda(b)x^{(\alpha-(p-1)\boldsymbol1)/p},
    &\alpha_i\equiv p-1\pmod p\text{ for all }i,\\
 0,&\text{otherwise}.
 \end{cases}
\]
This is additive and satisfies $T(a^pg)=aT(g)$. Since $H^{p-1}$
has degree $N(p-1)$, Proposition~\ref{prop:cyclic} implies
$T(H^{p-1})=1$. Thus
\begin{equation}\label{eq:splitting}
 \phi(g)=T(H^{p-1}g)
\end{equation}
is a Frobenius splitting. For each $j$,
\[
 H^{p-1}f_jg=f_j^p\left(\prod_{i\ne j}f_i^{p-1}\right)g,
\]
so $\phi$ preserves $(f_j)$. It therefore preserves the ideal
generated by any subset of the maximal permanents. The induced
quotient splitting proves reducedness: if $a^p$ is zero in the
quotient, applying the splitting gives $a=0$, and iteration eliminates
every nilpotent. The same argument works after arbitrary field
extension, so these quotients are geometrically reduced.

\subsection{Homogeneous specialization}
The next elementary fact will transfer both complete-intersection
and critical-locus bounds from one characteristic to zero.

\begin{lemma}\label{lem:specialization}
Let $J\subseteq\Z[x_1,\ldots,x_N]$ be generated by homogeneous
polynomials, and denote their generated ideals over $\Q$ and
$\F_p$ by $J_{\Q}$ and $J_p$. In every degree $d$,
\[
 \dim_{\Q}(\Q[x]/J_{\Q})_d
 \le\dim_{\F_p}(\F_p[x]/J_p)_d.
\]
In particular, the Krull dimension of the characteristic-zero
quotient is at most that of the characteristic-$p$ quotient.
\end{lemma}
\begin{proof}
For a fixed $d$, form the integer matrix whose columns are the
coefficient vectors of all products of a generator of $J$ with a
monomial of complementary degree. Its images over the two fields
are precisely the degree-$d$ parts of the generated ideals. Its rank
over $\Q$ is at least its rank modulo $p$, proving the Hilbert-function
inequality. The dimension assertion follows from the eventual
polynomial growth of the cumulative Hilbert function. This also
handles zero-dimensional graded quotients.
\end{proof}

\begin{proposition}\label{prop:lift}
Let $f_1,\ldots,f_s$ be positive-degree homogeneous polynomials over
$\Z$. Suppose their reductions modulo $p$ form a regular sequence
and generate a radical ideal. Then their characteristic-zero ideal
is a geometrically reduced complete intersection of height $s$.
\end{proposition}
\begin{proof}
Write $N$ for the number of variables and $d_i=\deg f_i$.
Lemma~\ref{lem:specialization} bounds the generic-fiber dimension
by $N-s$. The height theorem gives the reverse inequality, so the
characteristic-zero generators form a regular sequence. Both
quotients have Hilbert series
$\prod_i(1-z^{d_i})/(1-z)^N$.

Set $V=\Z_{(p)}$ and $B=V[x]/(f_1,\ldots,f_s)$.
Each graded piece $B_d$ is a finite $V$-module. Its rank is the
dimension of the corresponding generic-fiber piece; the dimension
of $B_d/pB_d$ is that of the special-fiber piece. The equality of
the Hilbert series makes these numbers equal. The structure theorem
for finite modules over the discrete valuation ring $V$ therefore
shows that each $B_d$ is free. In particular, $B$ is torsion-free.

Suppose $B\otimes_V\Q$ had a nonzero homogeneous nilpotent.
Choose a representative in some $B_d$ by clearing a power of $p$,
then divide out its maximal common power of $p$ in a $V$-basis.
The resulting element $a\in B_d$ has nonzero reduction modulo $p$.
A power of $a$ vanishes in the generic fiber, and torsion-freeness
of the appropriate graded piece makes that power zero already in
$B$. This contradicts reducedness of $B/pB$. The nilradical of a
graded ring is homogeneous, so the generic fiber is reduced.
A reduced finite-type algebra over the perfect field $\Q$ is
geometrically reduced. Extension to any characteristic-zero field
finishes the proof.
\end{proof}

Apply Proposition~\ref{prop:lift} to the full set of maximal
permanents, or to any subset, using the splitting and height
arguments above. This proves the complete-intersection and
reducedness assertions of Theorem~\ref{thm:prime}, including
\eqref{eq:hilbert}.

\section{Critical loci of row expansions}
\subsection{The full permanent}
Append a row $u=(u_1,\ldots,u_p)$ to $A$. Expansion along that
row gives
\begin{equation}\label{eq:rowexpansion}
 \per_p(A,u)=\sum_{j=1}^{p}u_jf_j(A).
\end{equation}
Over an algebraically closed field of characteristic $p$, its
critical equations are
\begin{equation}\label{eq:critical}
 f_1(A)=\cdots=f_p(A)=0,\qquad J_f(A)^{\mathsf t}u=0.
\end{equation}
For the rank-$r$ stratum of the base zero set,
Proposition~\ref{prop:rank} gives dimension at most
$N-(2p-r)$. The fiber in the $u$-variables has dimension $p-r$.
Hence the corresponding part of the critical locus has dimension
at most
\[
 N-(2p-r)+(p-r)=p^2-2p.
\]
There are finitely many ranks, so this bounds the entire critical
locus. The subspace consisting of matrices with two prescribed
zero rows has dimension $p^2-2p$ and satisfies all the cofactor
equations. Thus the codimension is exactly $2p$.

The critical equations are homogeneous with integer coefficients.
Lemma~\ref{lem:specialization} gives the same lower bound on
codimension in characteristic zero, and the two-zero-row subspace
again gives equality. This argument covers $p=2$ as well: the
critical locus of $x_{11}x_{22}+x_{12}x_{21}$ is the origin.

The permanent is a nonzero squarefree polynomial over every field.
Indeed, its degree in each variable is one; a repeated nonconstant
factor would depend on some variable and give degree at least two
in that variable. The hypersurface Jacobian criterion therefore
identifies its singular set with its gradient zero set after
extension to an algebraic closure. Equation~\eqref{eq:rowexpansion}
places the gradient zero set on the hypersurface even when the
characteristic divides the degree. This completes the proof of
\eqref{eq:primecodim} and Theorem~\ref{thm:prime}.

\subsection{Partial row expansions}
Fix $S\subseteq\{1,\ldots,p\}$ with $|S|=s\ge1$ and retain only
the variables $u_j$ with $j\in S$. The critical equations of $g_S$
are $f_j(A)=0$ for $j\in S$ and $J_{f_S}(A)^{\mathsf t}u=0$.
The subset version of Proposition~\ref{prop:rank} gives, at
Jacobian rank $r$, base dimension at most $N-(2s-r)$ and fiber
dimension $s-r$. Thus
\[
 \dim\Crit(g_S)\le N-s,
\]
or codimension at least $2s$ in $\A^{N+s}$.
Conversely the locus $u=0$, $f_j(A)=0$ for $j\in S$, has dimension
$N-s$, by the height statement. It lies in the critical locus.
Equality follows in characteristic $p$.

Lemma~\ref{lem:specialization} transfers the lower codimension
bound to characteristic zero. Proposition~\ref{prop:lift} gives
height $s$ for the subset ideal in that characteristic, so the
same $u=0$ locus attains equality there. The strength assertion
will follow in Section~\ref{sec:strength}.

\section{Propagation to arbitrary matrix sizes}
We require a dimension comparison for polynomial equations whose
highest-degree parts are known.

\begin{lemma}\label{lem:leading}
Let $h_1,\ldots,h_t$ be nonconstant polynomials over an algebraically
closed field, and let $h_i^{\mathrm{top}}$ be their highest
total-degree homogeneous parts. Then
\begin{equation}\label{eq:leading}
 \dim V(h_1,\ldots,h_t)
 \le\dim V(h_1^{\mathrm{top}},\ldots,h_t^{\mathrm{top}}).
\end{equation}
\end{lemma}
\begin{proof}
Let $J=(h_1,\ldots,h_t)$, and give $k[x]/J$ the total-degree
filtration. The associated graded quotient is
$k[x]/\operatorname{in}_{\deg}(J)$, where the initial ideal is
generated by the highest-degree parts of all elements of $J$.
The filtered algebra and its associated graded algebra have the
same Krull dimension, as follows from the growth of the degree
filtration. Since
\[
 (h_1^{\mathrm{top}},\ldots,h_t^{\mathrm{top}})
   \subseteq\operatorname{in}_{\deg}(J),
\]
the dimension inequality follows. The initial ideal may strictly
contain the ideal of the displayed leading parts; the inclusion in
this direction is all that is needed.
\end{proof}

\begin{proposition}\label{prop:border}
Over any field, the geometric codimensions satisfy
\begin{align}
 c_k(r,n)&\ge c_k(r-1,n-1)&& (2\le r\le n),\label{eq:borderrect}\\
 s_k(n)&\ge s_k(n-1)&& (n\ge3).\label{eq:bordersquare}
\end{align}
\end{proposition}
\begin{proof}
Work over an algebraic closure. Cover the nonzero matrices by the
finitely many open subsets on which a specified entry is nonzero.
By row and column permutations, it suffices to work with the entry
in position $(1,1)$. Fix the first row and the first column as
parameters, including $a=x_{11}\ne0$, and denote the free lower
block by $Y$.

First consider an $r\times n$ matrix. For each set $J$ of $r-1$
columns of $Y$, the maximal permanent using these columns and
column one has the form
\begin{equation}\label{eq:leadingrect}
 a\per(Y_J)+\text{terms of degree at most }r-2\text{ in }Y.
\end{equation}
To obtain $r-1$ free entries, all bottom rows must use free
columns, so the top row must use column one. This proves the
leading-term assertion. On every fixed-border fiber,
Lemma~\ref{lem:leading} bounds the dimension of the common zero
set of these selected equations by
\[
 (r-1)(n-1)-c_k(r-1,n-1).
\]
The full permanental zero set in the fiber is a subset of this
zero set. There are $r+n-1$ border parameters. The dimension of
the full zero set on the open chart is therefore at most
$rn-c_k(r-1,n-1)$. This proves \eqref{eq:borderrect} on each
nonzero chart. The omitted origin has dimension zero and does
not enlarge this bound.

For square critical loci, write the fixed-border matrix as
\[
 X=\begin{pmatrix}a&b^{\mathsf t}\\ c&Y\end{pmatrix},
 \qquad a\ne0,
\]
where $Y$ has row and column indices $1,\ldots,n-1$. Among the
gradient equations choose
$C_{ij}=\per(X_{\widehat{i+1},\widehat{j+1}})$ for all
$1\le i,j\le n-1$. Expanding according to the entries selected
in the first row and column gives
\begin{equation}\label{eq:leadingcofactor}
 C_{ij}=a\per(Y_{\widehat i,\widehat j})
 +\sum_{\substack{u\ne i\\v\ne j}}c_u b_v
       \per(Y_{\widehat{\{i,u\}},\widehat{\{j,v\}}}).
\end{equation}
Here a hat indicates deletion of the specified rows or columns,
and the permanent of the empty matrix is one. Each matching in
$C_{ij}$ has $n-1$ entries. If it uses $a$, the remaining $n-2$
entries lie in $Y$. If it avoids $a$, it must use two distinct
border entries, one in the first row and one in the first column;
only $n-3$ entries then lie in $Y$. Since the border is fixed,
the summands in the second term of \eqref{eq:leadingcofactor}
have degree at most $n-3$ in $Y$. Consequently
\[
 C_{ij}^{\mathrm{top}}
   =a\per(Y_{\widehat i,\widehat j}).
\]
These are $a$ times all the gradient coordinates of $\per_{n-1}(Y)$.
Lemma~\ref{lem:leading} bounds the common zero set of the selected
equations in each fixed-border fiber by
$(n-1)^2-s_k(n-1)$. The full critical fiber is a subset of that
zero set. Adding the $2n-1$ border parameters gives
$n^2-s_k(n-1)$. The nonzero charts and the origin cover the
critical locus, proving \eqref{eq:bordersquare}.
\end{proof}

\begin{proof}[Proof of Theorem~\ref{thm:alln}]
Iterate \eqref{eq:bordersquare} from $n$ down to $q$ and apply
Theorem~\ref{thm:prime} in characteristic zero. This gives
$s_k(n)\ge2q$. Two prescribed zero rows give the upper bound
$s_k(n)\le2n$ for every $n\ge2$.

Iterating \eqref{eq:borderrect} gives
$c_k(n-1,n)\ge c_k(q-1,q)=q$. There are $n$ maximal permanent
generators, so the height theorem gives $c_k(n-1,n)\le n$.
Finally, the prime number theorem implies $q/n\to1$. The two
limits follow by squeezing between the displayed bounds.
\end{proof}

\begin{remark}
The same propagation proves $s_k(n)\ge2p$ for every $n\ge p$
in a fixed positive characteristic $p$. It does not justify
using the largest prime below $n$ in that fixed characteristic:
the coefficient proof at order $q$ uses characteristic $q$.
\end{remark}

\section{Strength and scope}\label{sec:strength}
For clarity, all strength statements use the following convention.

\begin{definition}
For a nonzero homogeneous polynomial $f$ of degree at least two,
$\str(f)$ is the least $t$ such that
\[
 f=\sum_{i=1}^{t}a_ib_i,
\]
where $a_i,b_i$ are positive-degree homogeneous polynomials and
$\deg a_i+\deg b_i=\deg f$ for each $i$.
\end{definition}

\begin{lemma}\label{lem:strength}
Over any field,
\[
 \codim\Crit(f)\le2\str(f).
\]
\end{lemma}
\begin{proof}
For a decomposition into $t$ products, every first derivative
vanishes wherever all $2t$ factors vanish. Their common zero set
is nonempty, since the factors are homogeneous of positive degree,
and has codimension at most $2t$ by the height theorem. The
critical locus contains it, proving the claim.
\end{proof}

Expansion along one row expresses $\per_n$ as a sum of $n$
products, so its strength is at most $n$. Theorem~\ref{thm:alln}
and Lemma~\ref{lem:strength} give the lower bound $q$ in
characteristic zero. At prime order, Theorem~\ref{thm:prime}
gives exact strength $p$ in characteristic $p$ and zero.
Similarly, $g_S$ is a sum of $s$ displayed products and has
critical codimension $2s$, so its strength is $s$. This completes
Theorem~\ref{thm:partial}.

These results determine dimensions, rather than irreducible
components. They also do not determine the exact singular
codimension $2n$ at every composite size, or reducedness of the
corresponding maximal-permanent ideal at every composite number
of columns. Both questions remain natural extensions of the
prime-size argument.

The bounds concern specific algebraic invariants. Critical
codimension $2n$, by itself, does not imply a large arithmetic
circuit. For example,
\[
 h(z)=\sum_{i=1}^{n}\prod_{j=1}^{n}z_{ij}
\]
has a circuit with $O(n^2)$ arithmetic operations. Its critical
locus requires at least two zero entries in each row and hence
has codimension $2n$. Thus a circuit application of the
permanental geometry would require an additional argument.

\section{Exact finite checks}
The supplementary verification program uses integer arithmetic and
finite-field arithmetic. It checks the trace coefficient in the
small cases $p=2,3$, verifies instances of the generator-change
congruence in a truncated polynomial ring, and compares the
row-expansion critical equations with the full cofactor equations
on every $2\times2$ matrix over $\F_2$ and every $3\times3$
matrix over $\F_3$. A separate finite-field calculation checks
the cyclic fixed-point counts for additional primes. The program also
compares complete monomial expansions on both sides of
\eqref{eq:leadingcofactor} for every selected cofactor at sizes
$3\le n\le7$, including their degrees in the free block. These checks
are reproducible controls on the formulas. The all-prime assertions,
the geometric rank-stratum estimate, and the specialization and
border arguments are established by the proofs above.

\paragraph{Acknowledgment.}
The author acknowledges OpenAI GPT-6 for assistance with mathematical
exploration, exact computations, proof organization, and manuscript
preparation.

\begin{thebibliography}{3}\raggedright
\bibitem{BCMV}
A.~Boralevi, E.~Carlini, M.~Micha\l{}ek, and E.~Ventura,
\emph{On the codimension of permanental varieties},
Advances in Mathematics \textbf{461} (2025), article 110079.
\href{https://doi.org/10.1016/j.aim.2024.110079}{doi:10.1016/j.aim.2024.110079}.

\bibitem{Chau}
T.~Chau, \emph{The $F$-singularities of algebras defined by permanents},
Collectanea Mathematica \textbf{77} (2026), 843--859.
\href{https://doi.org/10.1007/s13348-025-00494-8}{doi:10.1007/s13348-025-00494-8}.

\bibitem{Knutson}
A.~Knutson, \emph{Frobenius splitting, point-counting, and degeneration},
preprint (2009).
\href{https://arxiv.org/abs/0911.4941v1}{arXiv:0911.4941v1}.
\end{thebibliography}
\end{document}
